\chapter{Stage 1: $QK^T$ GEMM Engine and Design Space Exploration}
\label{chap:stage1_gemm}

\section{Mathematical Formulation of Quantized Dot-Product}
Stage 1 computes the dot-product attention score between Query token $\mathbf{q}_i \in \text{INT8}^{1 \times d_k}$ and Key matrix $K \in \text{INT8}^{N \times d_k}$:
\begin{equation}
    s_{\text{raw}}[i, j] = \sum_{k=0}^{d_k-1} q_i[k] \cdot K[j, k]
\end{equation}
Where $q_i[k], K[j, k] \in [-128, 127]$. Each partial product requires 16 bits of dynamic range. Summing across $d_k = 16$ elements yields a maximum magnitude of:
\begin{equation}
    |s_{\text{raw}}| \le 16 \times (128 \times 128) = 262,144 < 2^{19}
\end{equation}
An internal 32-bit signed accumulator (\texttt{INT32}) is used to ensure absolute headroom against overflow across all token lengths.

\section{Architectural Variant 1A: Pure Logic Fabric Streaming Engine}
Variant 1A targets environments where all DSP blocks are occupied by convolutional backbones.
\begin{itemize}
    \item \textbf{Multiplier Implementation:} Uses 16 independent combinational multiplier macros synthesized entirely into CLB LUT6 logic fabric (\texttt{lut\_multiplier\_8x8.sv}).
    \item \textbf{Reduction Tree:} A 4-stage pipelined binary adder tree reduces 16 partial products to one INT32 sum.
    \item \textbf{Characteristics:} Consumes 0 DSP48E2 slices, but requires $\approx 5,150$ CLB LUTs. It sustains an initiation interval of $II = 1$ score/cycle.
\end{itemize}

\section{Architectural Variant 2A: High-Throughput DSP-Packed Streaming Engine}
The AMD UltraScale+ DSP48E2 hard slice features an asymmetric $18 \times 27$-bit signed multiplier. Standard 8-bit multiplication utilizes only a fraction of this capability. Variant 2A introduces the \texttt{pe\_slice\_dsp\_packing} microarchitecture, packing two independent INT8 multiplications into a single DSP slice.

\subsection{Mathematical Derivation of DSP Packing}
Let $A = q_i[k]$ be the common Query element, and let $B = K[j, 2k]$ and $C = K[j, 2k+1]$ be two Key elements. We construct the inputs to Port A and Port B as:
\begin{equation}
    \text{Port A} = A, \quad \text{Port B} = (B \ll 18) + C
\end{equation}
When multiplied inside the DSP48E2 core:
\begin{equation}
    \text{DSP\_Product} = A \times (B \cdot 2^{18} + C) = (A \times B) \cdot 2^{18} + (A \times C)
\end{equation}

\subsection{Borrow-Bit Phenomenon and Hardware $+1$ Sign Correction}
Because operand $C$ is a signed two's complement integer, when $C < 0$, its sign extension propagates high ones across the upper 18 bits of Port B. This behaves as an arithmetic borrow bit from the upper product:
\begin{equation}
    \text{If } A \times C < 0 \implies (\text{DSP\_Product} \ggg 18) = (A \times B) - 1
\end{equation}
To restore bit-exact fidelity, a hardware sign correction circuit inspects bit 15 of the lower product:
\begin{lstlisting}[language=Verilog, caption={Hardware Sign Correction in SystemVerilog}]
assign result_ac = dsp_product[15:0];
assign result_ab = (dsp_product >>> 18) + (dsp_product[15] ? 16'sd1 : 16'sd0);
\end{lstlisting}

\subsection{Numerical Demonstration}
Consider $A = 10$, $B = 20$, $C = -15$:
\begin{itemize}
    \item Expected: $A \times B = +200$, $A \times C = -150$.
    \item Port B: $(20 \ll 18) + (-15) = 20 \times 262,144 - 15 = 5,242,865$.
    \item DSP Product: $10 \times 5,242,865 = 52,428,650 = 199 \times 262,144 + 65,386$.
    \item Lower 16 bits: $65,386 \equiv -150 \pmod{65,536}$ (Bit-exact for $A \times C$).
    \item Bit 15 is 1 (negative result).
    \item Shifted upper part: $52,428,650 \gg 18 = 199$.
    \item Sign correction: $199 + 1 = \mathbf{200}$ (Bit-exact for $A \times B$).
\end{itemize}

Using this method, only \textbf{8 DSP48E2 slices} are required to compute the entire 16-element dot product, cutting LUT logic consumption by $47.6\%$.

\section{Architectural Variant 3B: 2D Spatial Systolic Array}
Variant 3B evaluates a $4 \times 4$ PE Output-Stationary Systolic Array (16 DSPs).
\begin{itemize}
    \item \textbf{Drain Overhead:} For small edge dimensions ($d_k = 16$), the computation phase requires 24 cycles, while draining partial sums requires 32 cycles. The array is idle for:
    \begin{equation}
        \text{Idle Stall Ratio} = \frac{32}{24 + 32} = 57.1\%
    \end{equation}
    Throughput drops to $32 / 56 = 0.57\text{ score/cycle}$.
    \item \textbf{Matrix Reordering Buffer:} Because systolic arrays output data in $4 \times 8$ matrix tiles, an auxiliary reordering buffer is required to convert tile data to sequential rows for Softmax.
\end{itemize}

\section{Comparative Evaluation of the Three Variants}
\begin{table}[htbp]
    \centering
    \caption{Synthesis and Performance Comparison Across the 3 Hardware Variants}
    \label{tab:gemm_variants_comparison}
    \begin{tabularx}{\textwidth}{l c c c c}
        \toprule
        \textbf{Metric} & \textbf{Variant 1A} & \textbf{Variant 2A} & \textbf{Variant 3B} & \textbf{Available (KV260)} \\
        \midrule
        DSP48E2 Slices          & 0 (0.0\%)      & \textbf{8 (0.64\%)}  & 16 (1.28\%)  & 1,248 \\
        CLB LUT6 Fabric         & 5,150 (4.4\%)  & \textbf{2,700 (2.3\%)}& 4,200 (3.6\%)& 117,120 \\
        Max Frequency ($F_{\max}$) & $\approx 330$~MHz & \textbf{$\approx 408$~MHz} & $\approx 250$~MHz & Target: 200~MHz \\
        Timing WNS (@ 200~MHz)  & $+1.98$~ns     & \textbf{$+0.051$~ns} & $+1.02$~ns   & $> 0$~ns \\
        Sustained Throughput    & 1.0 score/cyc  & \textbf{1.0 score/cyc} & 0.57 score/cyc & — \\
        Full Head Latency       & 1.048~ms       & \textbf{1.050~ms}    & 1.686~ms     & — \\
        Frame Rate (4 Heads)    & 238.5 FPS      & \textbf{238.2 FPS}   & 148.3 FPS    & Real-time (> 30) \\
        \bottomrule
    \end{tabularx}
\end{table}

Variant 2A achieves the ideal balance of minimal resource footprint, timing closure, and maximum real-time frame rate.

